\RequirePackage{iftex}
\ifPDFTeX
\documentclass[a4paper,12pt]{article}
\usepackage[margin=22mm]{geometry}
\usepackage[T1]{fontenc}
\usepackage[utf8]{inputenc}
\usepackage{graphicx}
\usepackage{CJKutf8}
\else
\documentclass[a4paper,12pt]{jsarticle}
\usepackage[dvipdfmx,margin=22mm]{geometry}
\usepackage[dvipdfmx]{graphicx}
\fi
\usepackage{amsmath,amssymb,amsthm}
\usepackage{enumitem}

\setlength{\parindent}{0pt}
\setlength{\parskip}{0.35em}
\setlength{\fboxsep}{8pt}
\setlist[itemize]{leftmargin=2em,itemsep=0.25em,topsep=0.35em}
\setlist[enumerate]{leftmargin=2.2em,itemsep=0.45em,topsep=0.35em}

\newcommand{\pointbox}[1]{%
\begin{center}
\fbox{\begin{minipage}{0.91\linewidth}#1\end{minipage}}
\end{center}
}

\begin{document}
\ifPDFTeX
\begin{CJK}{UTF8}{min}
\fi

\theoremstyle{definition}
\newtheorem*{definition}{定義}
\newtheorem*{theorem}{定理}
\newtheorem*{example}{例}
\newtheorem*{remark}{注意}
\renewcommand{\proofname}{\normalfont 証明}

\begin{center}
{\LARGE 数学1A 第13回レジュメ}\\[0.4em]
一様連続と平均値の定理
\end{center}

\section*{1. 一様連続性}

\begin{definition}
$f:D\to\mathbb{R}$ が一様連続であるとは、任意の $\varepsilon>0$ に対して
ある $\delta>0$ が存在して、任意の $x,y\in D$ について
\[
|x-y|<\delta\quad\Rightarrow\quad |f(x)-f(y)|<\varepsilon
\]
が成り立つことをいう。
\end{definition}

\begin{remark}
通常の連続性では $\delta$ は点に依存してよい。一様連続性では、同じ
$\delta$ を領域全体で使う。
\end{remark}

\begin{theorem}[Heine--Cantor の定理]
$f:[a,b]\to\mathbb{R}$ が連続ならば、$f$ は $[a,b]$ 上で一様連続である。
\end{theorem}

\begin{proof}
背理法で示す。一様連続でないとすると、ある $\varepsilon_0>0$ が存在して、
任意の $\delta>0$ に対して $x,y\in[a,b]$ があり、
\[
|x-y|<\delta,\qquad |f(x)-f(y)|\ge\varepsilon_0
\]
となる。$\delta=1/n$ として $x_n,y_n$ を取ると、$|x_n-y_n|<1/n$ である。

閉区間の性質より、$x_n$ のある部分列 $x_{n_k}$ は $x_0\in[a,b]$ に収束する。
$|x_{n_k}-y_{n_k}|\to0$ なので $y_{n_k}\to x_0$ でもある。連続性より
\[
f(x_{n_k})\to f(x_0),\qquad f(y_{n_k})\to f(x_0)
\]
となる。したがって $|f(x_{n_k})-f(y_{n_k})|\to0$ であり、
$|f(x_{n_k})-f(y_{n_k})|\ge\varepsilon_0$ に矛盾する。
\end{proof}

\section*{2. Rolle の定理}

\begin{theorem}[Rolle の定理]
$f:[a,b]\to\mathbb{R}$ が $[a,b]$ で連続、$(a,b)$ で微分可能で、
$f(a)=f(b)$ とする。このとき、ある $c\in(a,b)$ が存在して
\[
f'(c)=0
\]
となる。
\end{theorem}

\begin{proof}
最大値・最小値の定理より、$f$ は $[a,b]$ 上で最大値と最小値を取る。
もし最大値と最小値がともに端点で取られるなら、$f(a)=f(b)$ から $f$ は定数であり、
任意の $c\in(a,b)$ で $f'(c)=0$ である。

そうでない場合、最大値または最小値を取る点 $c$ が $(a,b)$ に存在する。内点で
極値を取る微分可能関数については $f'(c)=0$ である。
\end{proof}

\section*{3. 平均値の定理}

\begin{theorem}[平均値の定理]
$f:[a,b]\to\mathbb{R}$ が $[a,b]$ で連続、$(a,b)$ で微分可能であるとする。
このとき、ある $c\in(a,b)$ が存在して
\[
f'(c)=\frac{f(b)-f(a)}{b-a}
\]
となる。
\end{theorem}

\begin{proof}
\[
\varphi(x)=f(x)-\frac{f(b)-f(a)}{b-a}(x-a)
\]
とおく。$\varphi$ は $[a,b]$ で連続、$(a,b)$ で微分可能である。また
\[
\varphi(a)=f(a),\qquad
\varphi(b)=f(b)-\frac{f(b)-f(a)}{b-a}(b-a)=f(a)
\]
なので $\varphi(a)=\varphi(b)$ である。Rolle の定理より、ある $c\in(a,b)$ が
存在して $\varphi'(c)=0$ となる。よって
\[
0=f'(c)-\frac{f(b)-f(a)}{b-a}
\]
であり、主張が従う。
\end{proof}

\begin{example}
$f(x)=x^2$ を $[1,3]$ で考えると、
\[
\frac{f(3)-f(1)}{3-1}=\frac{9-1}{2}=4
\]
である。$f'(x)=2x$ なので、平均値の定理の点は $c=2$ である。
\end{example}

\pointbox{
平均値の定理は、区間全体の平均変化率が、どこか一点の瞬間変化率として
現れることを述べている。
}

\ifPDFTeX
\end{CJK}
\fi
\end{document}
